%% MODEL_replication_tidallanes.tex
%% Replication package introduction + Beijing tidal-lane counterfactual (Spec C)
%% 2026-04-29
%%
%% Compile: pdflatex MODEL_replication_tidallanes.tex  (twice for refs)
%% Requires: graphicx, booktabs, amsmath, hyperref, geometry

\documentclass[12pt, a4paper]{article}

\usepackage[utf8]{inputenc}
\usepackage[T1]{fontenc}
\usepackage[margin=2.5cm]{geometry}
\usepackage{amsmath, amssymb}
\usepackage{booktabs}
\usepackage{graphicx}
\usepackage{hyperref}
\usepackage{xcolor}
\usepackage{setspace}
\usepackage{parskip}
\usepackage{listings}
\usepackage{caption}
\usepackage{subcaption}
\usepackage{float}

\hypersetup{colorlinks=true, linkcolor=blue!60!black, citecolor=blue!60!black, urlcolor=teal}
\definecolor{codebg}{gray}{0.95}

\lstdefinestyle{pycode}{
  language=Python,
  backgroundcolor=\color{codebg},
  basicstyle=\ttfamily\footnotesize,
  keywordstyle=\bfseries\color{blue!70!black},
  commentstyle=\itshape\color{gray},
  stringstyle=\color{green!50!black},
  showstringspaces=false,
  breaklines=true,
  frame=single,
  framerule=0.4pt,
  xleftmargin=1em,
  xrightmargin=1em,
}

\title{\textbf{Replication Package Notes \& Application}}

\author{Yige Duan \and Xiaoruo Feng \and Yizhen Gu}
\date{\today}


\begin{document}
\maketitle
\tableofcontents
\newpage

% ============================================================
\section{Replication Package: Framework and Architecture}
\label{sec:replication}
% ============================================================

\subsection{Conceptual Framework}

This note adapts the exact-hat algebra of Ahlfeldt \& Albers (2022,
henceforth AA) to a Beijing counterfactual with direction-specific
commuting costs.\footnote{%
  Ahlfeldt, G.M.\ \& Albers, N.\ (2022).
  \emph{Commuting, wages, and urban form.}
  Reproduced and extended for Beijing AM-peak data.}

The city is represented as $N$ grid cells.  Each cell can host residents
$l_{R,i}$ and workers $l_{F,i}$, and bilateral commuting between residence
$i$ and workplace $j$ depends on an iceberg commuting cost
$\bar{t}_{ij}$.  The Beijing application keeps the AA spatial-equilibrium
objects but lets the cost matrix be directed: $\bar{t}_{ij}$ and
$\bar{t}_{ji}$ may differ because morning road congestion is asymmetric.

\paragraph{Welfare closure.}
The model records aggregate welfare through the scalar $\chi$, a
cost-index closure for the commuting equilibrium.  A stripped-down version
of the object is:
\begin{equation}
  \chi^{-1} \;=\; \sum_{i,j}
    \Bigl(\frac{\bar{u}_i\,\bar{v}_j}{\bar{t}_{ij}}\Bigr)^{\!\theta}
  \label{eq:chi}
\end{equation}
where $\bar{u}_i$ is residential amenity, $\bar{v}_j$ is workplace
productivity, and $\theta$ is the Fr\'echet dispersion parameter.  In the
implemented system, this scalar is solved jointly with residential and
workplace location shares.  Lower $\chi$ corresponds to higher expected
utility.

\paragraph{Exact-hat algebra.}
The counterfactual does not require levels of the unobserved fundamentals.
It solves for ratios of counterfactual to baseline values.  For any object
$z$, $\hat{z}=z^{\text{new}}/z^{\text{baseline}}$.  The welfare ratio is:
\begin{equation}
  \hat{\chi} = \frac{\chi^{\text{new}}}{\chi^{\text{baseline}}}
  \quad \Rightarrow \quad
  \hat{\chi} < 1 \;\text{(welfare \textbf{gain})}, \quad
  \hat{\chi} > 1 \;\text{(welfare \textbf{worsens})}.
  \label{eq:chihat}
\end{equation}

\subsection{Model Parameters}

\begin{table}[H]
\centering
\caption{Structural parameters}
\label{tab:params}
\begin{tabular}{lll}
\toprule
Parameter & Symbol & Value \\
\midrule
Fr\'echet dispersion         & $\theta$    & 6.83 \\
BPR congestion elasticity  & $\delta_1$  & 0.488 \\
Commuting--productivity coupling & $\lambda = \delta_1/\theta$ & 0.0714 \\
Residential agglomeration  & $\alpha$    & $-0.12$ \\
Workplace agglomeration    & $\beta$     & $-0.10$ \\
\bottomrule
\end{tabular}
\end{table}

The combined exponent used in hat-algebra inversion is:
\begin{equation}
  \frac{1 + \theta\lambda}{\theta} \;=\; \frac{1}{\theta} + \lambda \;\approx\; 0.218.
  \label{eq:exponent}
\end{equation}

\subsection{Input and Output Objects}

The AA solver requires three sparse input tables, each indexed by
consecutive 1-based node IDs:

\begin{table}[H]
\centering
\caption{Input tables for the AA commuting solver}
\label{tab:inputs}
\begin{tabular}{lll}
\toprule
Table & Key columns & Economic object \\
\midrule
Sparse adjacency
  & \texttt{i, j, Xi\_ij}
  & Baseline edge traffic input $\Xi_{ij}$ \\
Node masses
  & \texttt{node, l\_R, l\_F}
  & Baseline residents $l_{R,i}$ and workers $l_{F,j}$ per node \\
Commute OD
  & \texttt{i, j, L\_ij}
  & Observed road-commuting OD flows $L_{ij}$ \\
\bottomrule
\end{tabular}
\end{table}

The inherited AA file names are retained only as an input contract.  In this
Beijing application, the sparse-adjacency third column is not an iceberg
cost.  It is the baseline traffic-state input $\Xi_{ij}$, constructed from
the same raw-road path's observed travel time, free-flow travel time, and
lane proxy:
\begin{equation}
  \Xi_{ij} \;=\; \text{lanes}_{ij}
    \cdot \left(\frac{\max(\tau_{ij}^{\text{obs}},\,\tau_{ij}^{\text{ff}})}
                     {\tau_{ij}^{\text{ff}}}\right)^{1/\delta_1}.
  \label{eq:bpr_inversion}
\end{equation}
The output objects are the baseline equilibrium $(\chi,l_R,l_F)$ and, after
a policy shock, the counterfactual hats
$(\hat{\chi},\hat{l}_R,\hat{l}_F)$.  The OD table $L_{ij}$ enters as a
baseline commuting-flow object; the counterfactual solution reallocates
residents and workers across nodes rather than solving a separate
road-network route assignment.

\subsection{Solver Architecture}

\paragraph{Baseline equilibrium.}
Given $(\Xi,L,l_R,l_F)$ and the structural parameters, the baseline problem
solves for $(\chi,l_R,l_F)$ through a nested fixed point.  The outer loop
searches over $\chi$.  Conditional on that value, the inner loop iterates on
transformed variables $(x_1,x_2)$ until the implied residential and workplace
shares satisfy the equilibrium restrictions.  The transformation matrix is:
\begin{equation}
  \mathbf{X} =
  \begin{pmatrix}
    1 - \beta\theta & \dfrac{\theta\lambda(1-\alpha\theta)}{1+\theta\lambda} \\[8pt]
    \dfrac{\theta\lambda(1-\beta\theta)}{1+\theta\lambda} & 1 - \alpha\theta
  \end{pmatrix}
\end{equation}

\paragraph{Counterfactual equilibrium.}
A policy is represented as a directed cost shock
$\hat{\bar{t}}_{ij}$.  The solver applies the same nested structure in
hat form:
\begin{itemize}
  \item Given a candidate $\hat{\chi}$, iterate
    $(\hat{x}_1,\hat{x}_2)$ to recover $(\hat{l}_R,\hat{l}_F)$.
  \item Update $\hat{\chi}$ so that the equilibrium restrictions hold under
    the shocked cost matrix.
  \item Repeat until
    $|\log\hat{\chi}^{\text{new}}-\log\hat{\chi}|<\text{tol}$.
\end{itemize}


% ============================================================
\section{Application: Beijing Data to AA Inputs}
\label{sec:application}
% ============================================================

\subsection{Data Source}

The Beijing input bundle separates spatial units from road-path costs.
The retained centerline-grid pipeline supplies square grid IDs, node-level
residents and jobs, and the reachable AM road-commuting OD sample.  The
raw-edge-projection pipeline supplies directed AM path costs, free-flow
path costs, lane proxies, and BPR traffic-state inputs.  The active bundle
contains $2{,}723$ nodes, $3{,}666$ finite directed edges, and $377{,}708$
origin--destination rows.

\subsection{Edge Attributes Used}

\begin{table}[H]
\centering
\caption{Edge-level columns from the raw edge-projection cost file}
\label{tab:edge_cols}
\small
\begin{tabular}{p{4.5cm}p{2.3cm}p{6.0cm}}
\toprule
Column & Units & Description \\
\midrule
\texttt{tau\_obs\_min}      & minutes & Raw edge-projection AM path time \\
\texttt{tau\_ff\_min}       & minutes & Same path recomputed with 22:00--05:00 speeds \\
\texttt{lane\_resolved\_for\_bpr} & lanes & Path lane estimate, with numeric fallback \\
\texttt{xi\_bpr\_edgeproj} & traffic index & BPR-inverted baseline traffic input \\
\texttt{route\_length\_m}& metres  & Pair-specific raw-road path length \\
\texttt{home\_grid}, \texttt{work\_grid} & -- & Origin and destination grid IDs \\
\bottomrule
\end{tabular}
\end{table}

\subsection{Node Mapping}

Grid IDs are remapped to the consecutive 1-based node indices expected by
the AA solver:
\begin{equation}
  \texttt{node\_i} \;\longrightarrow\; \texttt{node\_i\_aa\_1based}
  \quad (\in \{1, \ldots, N\}), \qquad N = 2{,}723 \text{ nodes.}
\end{equation}

The node file is not restricted to the largest connected component of the
raw edge-projection graph.  A grid enters the AA node file if it has positive
residents or jobs, appears in the positive reachable-AM OD sample, or appears
as an endpoint of a finite raw edge-projection link.  The retained edge file
then contains only finite raw edge-projection AM links.  In the current
square input, this yields $2{,}723$ nodes and $3{,}666$ directed edges.

\subsection{Observed Flow Reconstruction (BPR Inversion)}

For each directed AM edge $(i,j)$, the traffic-state input is inferred from
the observed and free-flow path times:
\begin{align}
  \tau_{ij}^{\text{obs}} &= \tau_{ij}^{\text{ff}}
    \cdot \left(\frac{\xi_{ij}}{\text{lanes}_{ij}}\right)^{\delta_1}
  \;\Longrightarrow\;
  \xi_{ij} = \text{lanes}_{ij}
    \cdot \left(\frac{\max(\tau_{ij}^{\text{obs}},\,\tau_{ij}^{\text{ff}})}
                     {\tau_{ij}^{\text{ff}}}\right)^{1/\delta_1}
  \label{eq:xi}
\end{align}
This $\xi_{ij}$ enters the solver as the baseline traffic-state matrix.  The
observed time, free-flow time, and lane proxy are all measured on the same
pair-specific raw-road path.  The lane term prioritises the path-based
length-weighted lane estimate, then the path-based bottleneck estimate, and
finally a numeric fallback of two lanes.

The exported object is:
\begin{equation}
  \Xi_{ij}^{\text{BJ}} =
  \text{lanes}_{ij}^{\text{resolved}}
  \cdot
  \left(\frac{\max(\tau_{ij}^{\text{obs}},\,\tau_{ij}^{\text{ff}})}
             {\tau_{ij}^{\text{ff}}}\right)^{1/\delta_1},
  \label{eq:xi_bj_export}
\end{equation}
for each directed AM edge retained in the square-grid network.  This value is
written as the third column of the sparse adjacency input.

\subsection{Current AA Input Bundle}

The current Spec C run deploys the v2 Beijing input bundle.  Its edge-cost
and traffic-state inputs are based on raw edge-projection path travel times
and the BPR inversion in \eqref{eq:xi}.  The bundle preserves the AA file
contract but changes the economic content from Seattle inputs to Beijing
grid nodes, Beijing OD flows, and Beijing directed road-path costs.


% ============================================================
\section{Counterfactual: Tidal Lane Policy}
\label{sec:counterfactual}
% ============================================================

\subsection{Counterfactual Definition}

Here we implement a one-lane directional-capacity reallocation.  For each treated
bidirectional road pair, the AM direction with the longer observed travel
time receives one lane; the opposite direction loses one lane.  The slow
direction is defined from observed directed travel times.

The counterfactual asks how aggregate welfare and the spatial allocation of
residents and workers change when the 100 most asymmetric eligible road
pairs receive this tidal-lane-style reallocation.

\subsection{Implementation}

Eligible pairs must appear in the AM raw-edge-projection data, have more
than 200 metres of path length in both directions, and retain at least one
lane in the fast direction after reallocation.  They are ranked by:
\begin{equation}
  r_{ij} = \frac{\max(\tau_{ij}^{\text{obs}},\, \tau_{ji}^{\text{obs}})}
                {\min(\tau_{ij}^{\text{obs}},\, \tau_{ji}^{\text{obs}})}
  \label{eq:asymratio}
\end{equation}
The retained sample contains 100 treated pairs.  Their asymmetry ratios run
from 3.24 to 25.17.  Figure~\ref{fig:treated_grid_links} maps these treated
grid-pair links in the square-grid system.

Let $s$ denote the slow direction and $f$ the fast direction for a given pair.
The lane transfer is evaluated with the same raw-path observed time,
free-flow time, and lane proxy used in the v2 Beijing input bundle:
\begin{align}
  \tau_s^{\text{new}}
    &= \max\!\left[
        \tau_s^{\text{ff}} \cdot
        \left(\frac{\xi_s}{\text{lanes}_s + 1}\right)^{\!\delta_1},\;
        \tau_s^{\text{ff}}
      \right]
  \label{eq:tau_slow_new} \\[6pt]
  \tau_f^{\text{new}}
    &= \max\!\left[
        \tau_f^{\text{ff}} \cdot
        \left(\frac{\xi_f}{\text{lanes}_f - 1}\right)^{\!\delta_1},\;
        \tau_f^{\text{ff}}
      \right]
  \label{eq:tau_fast_new}
\end{align}
The free-flow floor prevents mechanically predicted times below the same
path's free-flow time.  The resulting directed iceberg-cost shocks are:
\begin{align}
  \hat{\bar{t}}_{s}
    &= \frac{\tau_s^{\text{new}}}{\tau_s^{\text{obs}}}, \\[4pt]
  \hat{\bar{t}}_{f}
    &= \frac{\tau_f^{\text{new}}}{\tau_f^{\text{obs}}}.
  \label{eq:tbarhat}
\end{align}

The counterfactual changes only the directed cost matrix for the 200 treated
directions.  It does not change the grid system, the OD matrix, free-flow
path times, total city population, or non-treated links.  The BPR step holds
baseline traffic inputs $\xi_s$ and $\xi_f$ fixed when computing the local
travel-time shock; the AA solver then re-solves the spatial commuting
equilibrium under the shocked cost matrix.  It does not solve a separate
Wardrop route-assignment problem.

\subsection{Results}

The local cost shock has the intended sign.  Capacity gains reduce travel
times in the slow direction by 10.6 percent on average.  Capacity losses
raise travel times in the fast direction by 20.7 percent on average.

\begin{table}[H]
\centering
\caption{Summary of iceberg cost shocks (100 treated pairs)}
\label{tab:tbar}
\begin{tabular}{lccc}
\toprule
Direction & Range of $\hat{\bar{t}}$ & Avg.\ $\tau$ change & Interpretation \\
\midrule
Slow (gains capacity) & $[0.876,\; 1.006]$ & $-10.6\%$ & Cost mostly falls \\
Fast (loses capacity) & $[1.201,\; 1.402]$ & $+20.7\%$ & Cost rises \\
\bottomrule
\end{tabular}
\end{table}

Figure~\ref{fig:tau_obs} plots the observed travel-time asymmetry among the
100 treated pairs.

\begin{table}[H]
\centering
\caption{Counterfactual welfare outcome ($\Delta = 1$ lane, top-100)}
\label{tab:welfare}
\begin{tabular}{ll}
\toprule
Metric & Value \\
\midrule
$\hat{\chi}$ & $0.9998669$ \\
$(\hat{\chi} - 1) \times 10^6$ & $-133.1$ \\
Welfare direction & \textbf{GAIN} \\
\bottomrule
\end{tabular}
\end{table}

The welfare index falls slightly, so the counterfactual is welfare-improving
in the model's metric.  The magnitude is modest: $\hat{\chi}=0.9998669$,
equivalent to a gain of 133.1 parts per million relative to the baseline
cost index.  The sign comes from the asymmetry in baseline congestion: the
slow direction receives extra capacity where observed travel times are
substantially higher.

The location responses are redistributions rather than changes in total city
population.  The model centre is defined internally as the centroid of the
union of all active square-grid polygons, rather than an externally assigned
CBD coordinate.  Relative to this centre, the average resident location moves
only 0.9 metres inward, and the average workplace location moves 0.4 metres
outward.  The distance-gradient pattern is therefore modest.  Resident shares
rise slightly in the inner tercile and fall in the outer tercile; workplace
shares rise most in the middle tercile.  Figures~\ref{fig:resident_distance_scatter}
and~\ref{fig:worker_distance_scatter} report the underlying node-level people
changes against distance to the model centre.


% ============================================================
\section*{Appendix: Figures}
% ============================================================
\addcontentsline{toc}{section}{Appendix: Figures}

\begin{figure}[H]
  \centering
  \includegraphics[width=\textwidth]{../L3_figs/fig3_tau_obs.pdf}
  \caption{Observed AM-peak travel times for the 100 treated road pairs.
    Each point is a bidirectional pair; the slow direction (y-axis) exceeds
    the fast direction (x-axis). Colour encodes the asymmetry ratio
    $r_{ij}=\tau^{\max}/\tau^{\min}$, ranging from 3.24 to 25.17
    in the retained sample.}
  \label{fig:tau_obs}
\end{figure}

\begin{figure}[H]
  \centering
  \includegraphics[width=0.88\textwidth]{../L3_figs/fig4_treated_grid_links.pdf}
  \caption{Treated grid-pair links in the square-grid system.
    Each segment connects one treated bidirectional pair; the segment is drawn
    from the fast AM direction toward the slow AM direction. Colour encodes the
    observed AM asymmetry ratio used for treatment ranking.}
  \label{fig:treated_grid_links}
\end{figure}

\begin{figure}[H]
  \centering
  \includegraphics[width=0.86\textwidth]{../L3_figs/fig5_resident_change_by_distance.pdf}
  \caption{Node-level resident change by distance to the model centre.
    The vertical axis is the change in resident counts implied by
    $l_{R,i}^{\text{base}}(\hat{l}_{R,i}-1)$; the horizontal axis is the
    distance from each grid cell to the geometric centre of the active
    square-grid system.  The black line is a LOWESS nonparametric fit.}
  \label{fig:resident_distance_scatter}
\end{figure}

\begin{figure}[H]
  \centering
  \includegraphics[width=0.86\textwidth]{../L3_figs/fig6_worker_change_by_distance.pdf}
  \caption{Node-level worker change by distance to the model centre.
    The vertical axis is the change in worker counts implied by
    $l_{F,i}^{\text{base}}(\hat{l}_{F,i}-1)$; the horizontal axis is the
    distance from each grid cell to the geometric centre of the active
    square-grid system.  The black line is a LOWESS nonparametric fit.}
  \label{fig:worker_distance_scatter}
\end{figure}


% ============================================================
\section{Open Questions and Refinement Directions}
\label{sec:open_questions}
% ============================================================

The current implementation should be read as a transparent AA-style
exact-hat counterfactual, not as a fully estimated Beijing structural model.
It takes the Beijing grid system, OD flows, directed raw-path costs, lane
proxies, and BPR-inverted traffic inputs as given.  The solver then evaluates
how a directed cost shock changes aggregate welfare and the spatial allocation
of residents and workers.

The most useful extension is a richer estimation layer built from Beijing
origin--destination flows and repeated travel-time observations.  With
observed commute shares
\begin{equation}
  \pi_{ij} = \frac{F_{ij}}{\sum_k F_{ik}},
\end{equation}
variation in $\tau_{ij}$ can identify the Fr\'echet dispersion parameter
$\theta$, subject to the usual concern that travel times and commuting demand
are jointly determined.  Link-level speed-flow variation can be used to
discipline the BPR congestion elasticity $\delta_1$.

With $\theta$, $\delta_1$, and richer data on wages or rents, the next version
could recover structural fundamentals $(\bar{u}_i,\bar{v}_j)$ and move from
hat-algebra counterfactuals toward level counterfactuals.  A more ambitious
version would replace the fixed-flow BPR shock with a Wardrop-consistent
traffic assignment, jointly determining link traffic and travel times inside
the spatial-equilibrium iteration.  That step is computationally heavier and
should follow, rather than precede, validation of the raw-edge-projection cost
inputs.

\end{document}
